\chapter{Heart Attack}

\section*{Overview}
A heart attack, or myocardial infarction, occurs when blood flow decreases or stops to a part of the heart, causing damage to the heart muscle. This is usually caused by a blockage in the coronary arteries that supply blood to the heart. Heart attack is a medical emergency that requires immediate treatment to restore blood flow and minimize heart damage.

\section*{Symptoms}
Classic symptoms of a heart attack include:
\begin{itemize}
    \item Chest pain or discomfort (pressure, squeezing, fullness)
    \item Pain radiating to arm, shoulder, neck, jaw, or back
    \item Shortness of breath
    \item Cold sweat
    \item Nausea or vomiting
    \item Lightheadedness or dizziness
    \item Unexplained fatigue
\end{itemize}
Women may experience atypical symptoms including jaw pain, back pain, nausea, and unusual fatigue without classic chest pain.


\section*{Causes}
Most heart attacks are caused by:
\begin{itemize}
    \item \textbf{Coronary artery disease:} Buildup of plaque (atherosclerosis) in coronary arteries
    \item \textbf{Plaque rupture:} Sudden rupture of cholesterol plaque triggering blood clot formation
    \item \textbf{Coronary spasm:} Temporary narrowing of coronary artery
    \item Risk factors include: high blood pressure, high cholesterol, smoking, diabetes, obesity, sedentary lifestyle, family history, age (>45 men, >55 women)
\end{itemize}

\section*{Diagnosis}
Heart attack diagnosis involves:
\begin{itemize}
    \item \textbf{Electrocardiogram (ECG/EKG):} Shows ST-segment changes indicating heart damage
    \item \textbf{Blood tests:} Cardiac enzymes (troponin, CK-MB) elevated when heart muscle damaged
    \item \textbf{Echocardiogram:} Ultrasound showing heart wall motion abnormalities
    \item \textbf{Coronary angiography:} Visualizes blockages in coronary arteries
    \item \textbf{Cardiac MRI:} Detailed imaging of heart damage
\end{itemize}
Time is critical - faster diagnosis leads to better outcomes.


\section*{Treatment}
Heart attack treatment focuses on restoring blood flow:
\begin{itemize}
    \item \textbf{Immediate:} Aspirin, oxygen, nitroglycerin, morphine for pain
    \item \textbf{Reperfusion therapy:}
        \begin{itemize}
            \item Percutaneous coronary intervention (PCI/stent) - preferred if available within 90 minutes
            \item Thrombolytic therapy ("clot-busting" drugs) if PCI unavailable
        \end{itemize}
    \item \textbf{Medications:} Beta-blockers, ACE inhibitors, statins, anticoagulants
    \item \textbf{Surgery:} Coronary artery bypass graft (CABG) for complex blockages
    \item \textbf{Cardiac rehabilitation:} Structured exercise and education program
\end{itemize}

\section*{Prevention}
Heart attack prevention includes:
\begin{itemize}
    \item Don't smoke and avoid secondhand smoke
    \item Control blood pressure (target <140/90 mmHg)
    \item Manage cholesterol (LDL <100 mg/dL, or <70 for high risk)
    \item Control diabetes (HbA1c <7\%)
    \item Maintain healthy weight (BMI 18.5-24.9)
    \item Exercise regularly (150 minutes moderate activity weekly)
    \item Eat heart-healthy diet (Mediterranean diet recommended)
    \item Limit alcohol consumption
    \item Manage stress
    \item Regular health screenings
\end{itemize}

\section*{When to Seek Medical Care}
Call emergency services immediately if you experience:
\begin{itemize}
    \item Chest pain or pressure lasting more than a few minutes
    \item Pain spreading to shoulder, arm, neck, jaw, or back
    \item Shortness of breath with or without chest discomfort
    \item Cold sweat, nausea, or lightheadedness
    \item Unexplained extreme fatigue
\end{itemize}
Do not drive yourself to the hospital. Early treatment saves heart muscle and lives.

\clearpage
